\chapter{INTRODUÇÃO}

\orientacao{A introdução contextualiza o tema, apresenta o problema de pesquisa e conduz o leitor até os objetivos. Comece pelo contexto geral e vá estreitando até o problema específico do seu trabalho. Use citações para sustentar afirmações (ex.: \texttt{\textbackslash cite\{autorAno\}}).}

\placeholder{Parágrafo 1 — Contextualização: apresente o cenário geral do tema escolhido. Por que esse assunto é relevante hoje? Inclua dados ou tendências recentes, se possível.}

\placeholder{Parágrafo 2 — Aprofundamento: detalhe aspectos específicos do contexto que se relacionam diretamente com o seu problema de pesquisa.}

\placeholder{Parágrafo 3 — Problema de pesquisa: formule a pergunta ou questão central que seu TCC pretende responder. Exemplo: ``Diante desse cenário, surge a questão: [sua pergunta de pesquisa]?''}

\placeholder{Parágrafo 4 — Delimitação: explique o recorte do estudo (o que será analisado, em que contexto, com quais limites).}

\section{Justificativa}

\orientacao{Explique por que vale a pena investigar esse tema. Aborde três dimensões quando possível: relevância acadêmica (lacuna na literatura), relevância social (impacto na sociedade) e relevância prática (aplicação ou utilidade).}

\placeholder{Justificativa acadêmica: qual lacuna ou debate científico seu trabalho contribui a preencher ou avançar?}

\placeholder{Justificativa social/prática: quem se beneficia dos resultados? Qual problema real ou institucional seu estudo ajuda a compreender ou resolver?}

\section{Objetivos}

\subsection{Objetivo Geral}

\orientacao{O objetivo geral deve ser um verbo no infinitivo que responda diretamente ao problema de pesquisa. Evite objetivos vagos como ``estudar'' ou ``analisar'' sem complemento. Prefira verbos como: investigar, comparar, propor, avaliar, identificar.}

\placeholder{Verbo no infinitivo + o quê + para quê/como. Exemplo: ``Analisar o desempenho de [X] em [contexto Y] para identificar [Z].''}

\subsection{Objetivos Específicos}

\orientacao{Liste de 3 a 6 objetivos específicos que, juntos, conduzem ao objetivo geral. Cada item deve ser mensurável ou verificável. Use o ambiente \texttt{itemize}.}

\begin{itemize}
    \item \placeholder{Objetivo específico 1 — ex.: revisar a literatura sobre o tema};
    \item \placeholder{Objetivo específico 2 — ex.: definir a metodologia de coleta/análise};
    \item \placeholder{Objetivo específico 3 — ex.: aplicar a técnica/método proposto};
    \item \placeholder{Objetivo específico 4 — ex.: interpretar os resultados esperados};
    \item \placeholder{Objetivo específico 5 — opcional}.
\end{itemize}
